\documentclass[11pt]{article}
\usepackage[margin=1in]{geometry}
\usepackage{amsmath,amssymb,amsthm}
\usepackage{enumitem}
\usepackage{booktabs}

\setlength{\parindent}{0pt}
\setlength{\parskip}{6pt}

\newcounter{qnum}
\newcommand{\question}[2]{\stepcounter{qnum}\subsection*{Question \theqnum \quad (#1 points) \hfill #2}}

\begin{document}

\begin{center}
{\LARGE\bfseries STAT GU4204 Section 001 --- Final Exam}\\[4pt]
{\large Variant E: Hybrid / Most-Likely}\\[6pt]
{\large Spring 2026 \hspace{1.5cm} Total: 100 points \hspace{1.5cm} Time: 150 minutes}\\[4pt]
\end{center}

\medskip

Name: \underline{\hspace{3in}} \hspace{1cm} UNI: \underline{\hspace{1.5in}}

\medskip

\textbf{Instructions.} This is an in-class, closed-book exam. You may bring \textbf{one} double-sided, hand-written 8.5$\times$11 cheat sheet. No calculators, computers, books, notes, or phones are permitted. Show all work to get full credit. Justify every step. Box your final answers.

\medskip

\textbf{Conventions.} Throughout the exam:
\begin{itemize}[leftmargin=2em,nosep]
    \item ``Random sample of size $n$ from $f(\cdot;\theta)$'' means $X_1,\ldots,X_n$ are i.i.d.\ with density (or pmf) $f(\cdot;\theta)$.
    \item Exponential$(\theta)$ has density $f(x;\theta) = \theta e^{-\theta x}$ on $x>0$ (\textbf{rate} parameterization), with $E[X]=1/\theta$ and $\mathrm{Var}(X)=1/\theta^2$.
    \item Beta$(\alpha,\beta)$ has density proportional to $p^{\alpha-1}(1-p)^{\beta-1}$ on $(0,1)$, with mean $\alpha/(\alpha+\beta)$.
    \item ``Regularity conditions'' are the standard (A.1)--(A.3) of DeGroot \& Schervish for the Cram\'er-Rao bound and for asymptotic normality of the MLE.
\end{itemize}

\medskip

\textbf{Quantile Table (use these as needed):}
\begin{center}\small
\begin{tabular}{ll}
\toprule
\textbf{Normal} & $z_{0.10}=1.282 \quad z_{0.05}=1.645 \quad z_{0.025}=1.960 \quad z_{0.01}=2.326 \quad z_{0.005}=2.576$ \\
\textbf{$t$} & $t_{0.025,7}=2.365 \quad t_{0.025,8}=2.306 \quad t_{0.025,9}=2.262 \quad t_{0.025,14}=2.145$ \\
            & $t_{0.025,15}=2.131 \quad t_{0.025,18}=2.101 \quad t_{0.025,19}=2.093 \quad t_{0.025,24}=2.064$ \\
            & $t_{0.05,8}=1.860 \quad t_{0.05,14}=1.761 \quad t_{0.05,15}=1.753 \quad t_{0.05,18}=1.734 \quad t_{0.05,19}=1.729$ \\
\textbf{$\chi^2$} & $\chi^2_{0.95,7}=14.07 \quad \chi^2_{0.95,8}=15.51 \quad \chi^2_{0.05,8}=2.733$ \\
                  & $\chi^2_{0.975,9}=19.02 \quad \chi^2_{0.025,9}=2.700$ \\
                  & $\chi^2_{0.95,1}=3.841 \quad \chi^2_{0.95,2}=5.991$ \\
\textbf{$F$} & $F_{0.05,4,4}=6.39 \quad F_{0.05,7,7}=3.79 \quad F_{0.025,7,7}=4.99 \quad F_{0.05,9,9}=3.18$ \\
             & $F_{0.95,9,9} = 1/F_{0.05,9,9} \approx 0.31$ \\
\textbf{Aux.} & $\Phi(0.855) \approx 0.80 \quad \Phi(1.645)=0.95 \quad \Phi(1.960)=0.975$ \\
\bottomrule
\end{tabular}
\end{center}

\bigskip
\hrule
\bigskip

% ====================================================================
% PART A
% ====================================================================
\section*{Part A: Short Conceptual Questions \hfill (10 points)}

\textit{Brief but complete answers. State theorems precisely and give 1--3 sentences of justification.}

\question{4}{Sufficiency and the Factorization Criterion}
Define what it means for a statistic $T = r(\mathbf{X})$ to be \emph{sufficient} for the parameter $\theta$. Then state the Fisher--Neyman Factorization Criterion. As a concrete example, let $X_1,\ldots,X_n$ be a random sample from $\mathrm{Bernoulli}(p)$ and use the criterion to identify a sufficient statistic for $p$. Show the factorization explicitly.

\question{3}{Cram\'er-Rao Lower Bound}
State the Cram\'er-Rao Lower Bound (CRLB) for the variance of an unbiased estimator $T$ of $\theta$ based on a random sample $X_1,\ldots,X_n$. Now consider $X_1,\ldots,X_n$ i.i.d.\ $\mathrm{Uniform}(0,\theta)$. Identify which one of the regularity conditions (A.1)--(A.3) is violated, and explain in one sentence why this prevents direct application of the CRLB.

\question{3}{Neyman-Pearson Lemma}
State the Neyman-Pearson Lemma for testing a simple null $H_0: \theta = \theta_0$ versus a simple alternative $H_1: \theta = \theta_1$. Be precise about the form of the optimal test and what optimality property it guarantees.

\bigskip
\hrule
\bigskip

% ====================================================================
% PART B
% ====================================================================
\section*{Part B: Mid-Length Problems \hfill (30 points)}

\question{15}{MLE, Sufficiency, and Fisher Information for the Exponential}
Let $X_1,\ldots,X_n$ be a random sample from the Exponential distribution with rate parameter $\theta > 0$, i.e.\ $f(x;\theta) = \theta e^{-\theta x}$ for $x > 0$.

\begin{enumerate}[label=(\alph*),leftmargin=2em]
    \item (5 pts) Find the maximum likelihood estimator $\hat\theta_n$ of $\theta$. Show all steps: write down the likelihood and log-likelihood, take the derivative, solve the score equation, and verify that you have a maximum.
    \item (4 pts) Using the Fisher--Neyman factorization criterion, show that $T = \sum_{i=1}^n X_i$ is sufficient for $\theta$. Identify the functions $u(\mathbf{x})$ and $\nu(t,\theta)$ explicitly.
    \item (4 pts) Compute the Fisher information $I_n(\theta)$ in the full sample. Is $\hat\theta_n$ an efficient estimator of $\theta$? Justify your answer carefully (in particular, comment on whether $\hat\theta_n$ is unbiased).
    \item (2 pts) State the asymptotic distribution of $\sqrt{n}(\hat\theta_n - \theta)$ as $n \to \infty$. Identify the asymptotic variance explicitly.
\end{enumerate}

\question{15}{Bayesian Inference for a Bernoulli Proportion}
Let $X_1,\ldots,X_n$ be a random sample from $\mathrm{Bernoulli}(p)$ with $p \in (0,1)$ unknown. Assume the prior distribution on $p$ is $\mathrm{Beta}(\alpha,\beta)$.

\begin{enumerate}[label=(\alph*),leftmargin=2em]
    \item (5 pts) Show that the posterior distribution of $p$ given $\mathbf{X} = (X_1,\ldots,X_n)$ is $\mathrm{Beta}\!\left(\alpha + \sum_{i=1}^n X_i,\; \beta + n - \sum_{i=1}^n X_i\right)$. Show all steps using the relation posterior $\propto$ likelihood $\times$ prior.
    \item (4 pts) Find the Bayes estimator $\hat p_B$ of $p$ under squared-error loss. Express your answer as a weighted average of the prior mean and the sample mean, and identify the weights.
    \item (3 pts) Suppose $\alpha = \beta = 1$ (uniform prior on $(0,1)$), $n = 10$, and $\sum_{i=1}^{10} X_i = 7$. Compute the numerical value of the Bayes estimator $\hat p_B$ and compare it to the maximum likelihood estimator $\hat p_n$.
    \item (3 pts) State (without proof) the asymptotic relationship between $\hat p_B$ and $\hat p_n$ as $n \to \infty$ with the prior $\mathrm{Beta}(\alpha,\beta)$ fixed. Briefly explain why this happens.
\end{enumerate}

\bigskip
\hrule
\bigskip

% ====================================================================
% PART C
% ====================================================================
\section*{Part C: Major Problems \hfill (60 points)}

\question{20}{$t$-Test, Confidence Interval, and Test/CI Duality}
Let $X_1,\ldots,X_n$ be a random sample from $N(\mu,\sigma^2)$ where both $\mu$ and $\sigma^2 > 0$ are unknown. Suppose $n = 16$, the observed sample mean is $\overline{X}_n = 12.5$, and the sample variance (with the unbiased $n-1$ denominator) is $s_n^2 = 4.0$.

\begin{enumerate}[label=(\alph*),leftmargin=2em]
    \item (8 pts) Test $H_0: \mu = 11$ versus $H_1: \mu \neq 11$ at significance level $\alpha_0 = 0.05$. Write down the test statistic
    \[ U_n = \frac{\sqrt{n}(\overline{X}_n - \mu_0)}{s_n}, \]
    state its distribution under $H_0$, give the critical value (use $t_{0.025,15} = 2.131$), and state your conclusion.
    \item (5 pts) Construct a 95\% confidence interval for $\mu$. Then verify explicitly that the test in part (a) is equivalent to checking whether $\mu_0 = 11$ lies inside this CI. State the duality theorem you are invoking.
    \item (4 pts) Now consider the one-sided problem $H_0: \mu \le 11$ versus $H_1: \mu > 11$ at level $\alpha_0 = 0.05$. Compute the new rejection threshold (use $t_{0.05,15} = 1.753$), give the value of the test statistic, and state your decision.
    \item (3 pts) For the two-sided test in part (a), let $\Lambda(\mathbf{X})$ be the likelihood ratio statistic. State the asymptotic null distribution of $-2 \log \Lambda(\mathbf{X})$ as $n \to \infty$, naming the theorem you invoke and identifying the degrees of freedom.
\end{enumerate}

\question{20}{Neyman-Pearson Test, Type I/II Errors, and UMP}
Let $X_1,\ldots,X_n$ be a random sample from $N(\theta, 1)$ with $n = 25$. Consider testing
\[ H_0: \theta = 0 \quad \text{vs.} \quad H_1: \theta = 0.5. \]

\begin{enumerate}[label=(\alph*),leftmargin=2em]
    \item (6 pts) Apply the Neyman-Pearson Lemma to derive the most powerful test. Show that the rejection region is of the form $\{\overline{X}_n > c\}$ for some constant $c$ that you should identify symbolically (in terms of $n$, $\theta_0=0$, $\theta_1=0.5$, and the threshold $k$ from the lemma).
    \item (4 pts) Find the value of $c$ such that the test has level $\alpha_0 = 0.05$. Use $z_{0.05} = 1.645$.
    \item (5 pts) Compute the Type I error probability $\alpha$ (and verify it equals $\alpha_0$). Then compute the Type II error probability $\beta = \mathbb{P}_{\theta = 0.5}(\overline{X}_n \le c)$. Express $\beta$ as $\Phi(\cdot)$ and use $\Phi(0.855) \approx 0.80$ to give a numerical value.
    \item (3 pts) Now consider the composite problem $H_0: \theta \le 0$ versus $H_1: \theta > 0$ at level $\alpha_0 = 0.05$. Argue that the same test you derived in (a)--(b) is uniformly most powerful (UMP) of level $\alpha_0$. You may invoke the monotone likelihood ratio (MLR) property if you state it.
    \item (2 pts) For the simple-vs-simple problem in (a), find the smallest sample size $n^*$ such that the most powerful level-0.05 test has power at least 0.95 at $\theta = 0.5$. (You may leave $n^*$ in symbolic form involving $z$-quantiles, then round up to the nearest integer.)
\end{enumerate}

\question{20}{Two-Sample $t$-Test, CI, and $F$-Test}
Two independent random samples are drawn from normal populations with a common (initially assumed) unknown variance $\sigma^2$:

\begin{itemize}[leftmargin=2em,nosep]
    \item Sample $\mathbf{X}$: $m = 10$ observations, $\overline{X}_m = 18.0$, $S_X^2 = \sum_{i=1}^m (X_i - \overline{X}_m)^2 = 36.0$.
    \item Sample $\mathbf{Y}$: $n = 10$ observations, $\overline{Y}_n = 15.5$, $S_Y^2 = \sum_{j=1}^n (Y_j - \overline{Y}_n)^2 = 28.0$.
\end{itemize}

\begin{enumerate}[label=(\alph*),leftmargin=2em]
    \item (3 pts) Compute the pooled variance estimate
    \[ s_p^2 = \frac{S_X^2 + S_Y^2}{m + n - 2}. \]
    \item (6 pts) Test $H_0: \mu_X = \mu_Y$ versus $H_1: \mu_X \neq \mu_Y$ at significance level $\alpha_0 = 0.05$. Compute the two-sample $t$-statistic
    \[ U = \frac{\overline{X}_m - \overline{Y}_n}{s_p\sqrt{\tfrac{1}{m} + \tfrac{1}{n}}}, \]
    state its distribution under $H_0$, give the critical value (use $t_{0.025,18} = 2.101$), and state your conclusion.
    \item (4 pts) Construct a 95\% confidence interval for the mean difference $\mu_X - \mu_Y$. Verify your CI is consistent with the test conclusion in (b).
    \item (5 pts) Now suppose we no longer assume equal variances. Test $H_0: \sigma_X^2 = \sigma_Y^2$ versus $H_1: \sigma_X^2 \neq \sigma_Y^2$ at significance level $\alpha_0 = 0.10$. Compute the $F$-statistic
    \[ V = \frac{S_X^2/(m-1)}{S_Y^2/(n-1)}, \]
    state its distribution under $H_0$, give the rejection region (use $F_{0.05,9,9} = 3.18$ and $F_{0.95,9,9} \approx 0.31$), and state your conclusion.
    \item (2 pts) For the test in part (b), state the asymptotic null distribution of $-2\log\Lambda$ where $\Lambda$ is the likelihood ratio statistic, and verify that the degrees of freedom match what Wilks' theorem predicts (count constraints).
\end{enumerate}

\bigskip
\hrule
\bigskip

\begin{center}
\textbf{End of exam.} \quad Total: 100 points.\\
Check that you answered all parts. Justification is required for full credit.
\end{center}

\end{document}
